% =============================================================================
\chapter{Thực nghiệm}\label{chap:thucnghiem}
% =============================================================================

\section{Mục tiêu và thiết lập}

Thực nghiệm trả lời câu hỏi của phần nâng cao: \emph{bậc $m$ ảnh hưởng thế nào tới số nút phải truy xuất khi tìm
kiếm, chiều cao cây và thời gian chạy}, với số mặt hàng tăng tới $10^5$.

\begin{itemize}
  \item \textbf{Dữ liệu:} $n \in \{10^3,\ 5\cdot10^3,\ 10^4,\ 2\cdot10^4,\ 5\cdot10^4,\ 10^5\}$ mã vạch 13 chữ số khác
        nhau; $m \in \{3, 5, 10\}$.
  \item \textbf{Cách đo} (hàm \code{doBTree}): xáo trộn thứ tự rồi chèn cả $n$ mặt hàng; xáo trộn lại rồi tìm lại
        cả $n$ mã. Thời gian đo bằng \code{std::chrono::steady\_clock}; số nút truy xuất trung bình
        $= \mathit{soNutTruyXuat} / n$ trong pha tìm. Bộ sinh số ngẫu nhiên xorshift32~\cite{xorshift} tự cài, hạt
        giống cố định, nên kết quả lặp lại được trên mọi máy.
  \item \textbf{Lặp lại:} mỗi cấu hình chạy 3 lần với 3 hạt giống khác nhau, lấy trung bình.
  \item \textbf{Máy chạy:} Intel Xeon 2,10\,GHz (4 nhân ảo, chương trình chạy đơn luồng), Linux, \code{g++} 13.3,
        \code{-O2}. Lệnh: \code{make bench}; kết quả ghi vào \code{ketqua/ketqua\_btree.csv}.
\end{itemize}

\section{Kết quả}

\pgfplotstableread[col sep=comma]{\CSV}\bangKQ

\begin{table}[htbp]
\centering\small
\caption{Kết quả thực nghiệm B-Tree (trung bình 3 lần chạy)}
\label{tab:thucnghiem}
\pgfplotstabletypeset[
  columns={n,m,chieuCao,soNut,tbNutTruyXuat,tgChenMs,tgTimMs},
  columns/n/.style={column name={$n$},int detect,1000 sep={\,}},
  columns/m/.style={column name={$m$},int detect},
  columns/chieuCao/.style={column name={Chiều cao},int detect},
  columns/soNut/.style={column name={Số nút},int detect,1000 sep={\,}},
  columns/tbNutTruyXuat/.style={column name={Nút truy xuất TB},fixed,fixed zerofill,precision=2,use comma},
  columns/tgChenMs/.style={column name={Chèn $n$ khóa (ms)},fixed,fixed zerofill,precision=2,use comma},
  columns/tgTimMs/.style={column name={Tìm $n$ khóa (ms)},fixed,fixed zerofill,precision=2,use comma},
  every head row/.style={before row=\toprule,after row=\midrule},
  every last row/.style={after row=\bottomrule},
  every nth row={3}{before row=\midrule},
]\bangKQ
\end{table}

\begin{figure}[htbp]
\centering
\begin{tikzpicture}
\begin{axis}[width=0.66\linewidth,height=7cm,xmode=log,log basis x=10,
  xlabel={Số mặt hàng $n$},ylabel={Số nút truy xuất trung bình mỗi lần tìm},
  ymin=0,ymax=17,grid=major,grid style={black!12},legend pos=outer north east,legend cell align=left,
  legend style={font=\footnotesize,align=left},tick label style={font=\footnotesize},
  /pgf/number format/use comma]
  \addplot[thick,blue!70!black,mark=*] table[x=n,y=tbNutTruyXuat,col sep=comma,discard if not={m}{3}]{\CSV};
  \addlegendentry{$m = 3$ (đo)}
  \addplot[thick,orange!85!black,mark=square*] table[x=n,y=tbNutTruyXuat,col sep=comma,discard if not={m}{5}]{\CSV};
  \addlegendentry{$m = 5$ (đo)}
  \addplot[thick,green!50!black,mark=triangle*] table[x=n,y=tbNutTruyXuat,col sep=comma,discard if not={m}{10}]{\CSV};
  \addlegendentry{$m = 10$ (đo)}
  \pgfplotsinvokeforeach{2,3,5}{%
    \addplot[dashed,black!45,forget plot,samples at={1000,1500,2000,3000,5000,7000,10000,15000,20000,30000,50000,70000,100000}]
      {1+ln((x+1)/2)/ln(#1)};}
  \addlegendimage{dashed,black!45}\addlegendentry{cận trên chiều cao\\$1+\log_t\frac{n+1}{2}$}
  \node[font=\scriptsize,black!55,anchor=west] at (axis cs:100000,16.6) {$t{=}2$};
  \node[font=\scriptsize,black!55,anchor=west] at (axis cs:100000,10.85) {$t{=}3$};
  \node[font=\scriptsize,black!55,anchor=west] at (axis cs:100000,7.7) {$t{=}5$};
\end{axis}
\end{tikzpicture}
\caption{Số nút truy xuất trung bình mỗi lần tìm theo $n$ và $m$; nét đứt là cận trên lý thuyết của chiều cao}
\label{fig:nuttruyxuat}
\end{figure}

\begin{figure}[htbp]
\centering
\begin{subfigure}[b]{0.49\linewidth}\centering
\begin{tikzpicture}
\begin{axis}[width=\linewidth,height=6cm,xmode=log,ymode=log,xlabel={$n$},ylabel={Thời gian (ms)},
  grid=major,grid style={black!12},legend pos=north west,legend style={font=\scriptsize},
  tick label style={font=\scriptsize},label style={font=\footnotesize}]
  \addplot[thick,blue!70!black,mark=*] table[x=n,y=tgChenMs,col sep=comma,discard if not={m}{3}]{\CSV}; \addlegendentry{$m=3$}
  \addplot[thick,orange!85!black,mark=square*] table[x=n,y=tgChenMs,col sep=comma,discard if not={m}{5}]{\CSV}; \addlegendentry{$m=5$}
  \addplot[thick,green!50!black,mark=triangle*] table[x=n,y=tgChenMs,col sep=comma,discard if not={m}{10}]{\CSV}; \addlegendentry{$m=10$}
\end{axis}
\end{tikzpicture}
\caption{Tổng thời gian chèn $n$ khóa}
\end{subfigure}\hfill
\begin{subfigure}[b]{0.49\linewidth}\centering
\begin{tikzpicture}
\begin{axis}[width=\linewidth,height=6cm,xmode=log,ymode=log,xlabel={$n$},ylabel={Thời gian (ms)},
  grid=major,grid style={black!12},legend pos=north west,legend style={font=\scriptsize},
  tick label style={font=\scriptsize},label style={font=\footnotesize}]
  \addplot[thick,blue!70!black,mark=*] table[x=n,y=tgTimMs,col sep=comma,discard if not={m}{3}]{\CSV}; \addlegendentry{$m=3$}
  \addplot[thick,orange!85!black,mark=square*] table[x=n,y=tgTimMs,col sep=comma,discard if not={m}{5}]{\CSV}; \addlegendentry{$m=5$}
  \addplot[thick,green!50!black,mark=triangle*] table[x=n,y=tgTimMs,col sep=comma,discard if not={m}{10}]{\CSV}; \addlegendentry{$m=10$}
\end{axis}
\end{tikzpicture}
\caption{Tổng thời gian tìm $n$ khóa}
\end{subfigure}
\caption{Thời gian chèn và tìm theo $n$ (trục log--log)}
\label{fig:thoigian}
\end{figure}

\section{Phân tích}

\paragraph{Số nút truy xuất giảm mạnh khi tăng $m$.} Với $n = 10^5$, mỗi lần tìm đi qua trung bình
13,25 nút khi $m = 3$, 8,63 nút khi $m = 5$ và chỉ 5,84 nút khi $m = 10$ -- giảm khoảng \textbf{2,3 lần} từ
$m = 3$ lên $m = 10$ (Hình~\ref{fig:nuttruyxuat}). Đường biểu diễn gần như thẳng trên trục $n$ dạng log, đúng với
$O(\log_t n)$: nhân $n$ lên 10 lần chỉ làm tăng thêm khoảng $\log_t 10$ nút ($\approx 3{,}3$ với $m=3$;
$\approx 1{,}4$ với $m = 10$). Cây nhị phân cân bằng (như AVL) với $n = 10^5$ cần trung bình xấp xỉ
$\log_2 n \approx 16{,}6$ nút mỗi lần tìm~\cite{knuth}, nhiều hơn B-Tree bậc 10 khoảng 2,8 lần.

\paragraph{Số nút trung bình luôn nhỏ hơn chiều cao.} Lần tìm kết thúc ngay khi gặp khóa, kể cả ở nút trong,
nên chỉ những khóa nằm ở lá mới cần đi hết $h$ nút. Với $n = 10^5$, tỉ lệ khóa nằm ở nút trong là 42,9\,\%
($m=3$), 27,1\,\% ($m=5$) và 14,1\,\% ($m=10$). Nhóm kiểm tra lại bằng cách đếm độ sâu từng khóa trong cây: độ sâu
trung bình bằng đúng 13,25; 8,63 và 5,84 -- trùng với số nút truy xuất đo được.

\paragraph{Chiều cao nằm đúng trong cận lý thuyết.} Với B-Tree $n$ khóa, chiều cao $h$ thỏa
$\lceil\log_m(n+1)\rceil \le h \le 1 + \log_t\frac{n+1}{2}$~\cite{clrs}. Bảng~\ref{tab:canchieucao} cho thấy mọi
chiều cao đo được nằm trong khoảng này và gần cận dưới hơn: cây thực tế thấp hơn trường hợp xấu nhất vì các nút
không chỉ chứa số khóa tối thiểu.

\begin{table}[htbp]
\centering\small
\caption{Chiều cao đo được so với cận lý thuyết và tỉ lệ lấp đầy nút}
\label{tab:canchieucao}
\begin{tabular}{@{}r c c c c c@{}}
\toprule
$n$ & $m$ & Cận dưới $\lceil\log_m(n+1)\rceil$ & \textbf{Chiều cao đo} & Cận trên $\lfloor 1+\log_t\frac{n+1}{2}\rfloor$ & Lấp đầy nút \\
\midrule
1\,000 & 3 & 7 & \textbf{8} & 9 & 68,1\,\% \\
1\,000 & 5 & 5 & \textbf{6} & 6 & 68,5\,\% \\
1\,000 & 10 & 4 & \textbf{4} & 4 & 67,8\,\% \\
\midrule
10\,000 & 3 & 9 & \textbf{11} & 13 & 67,1\,\% \\
10\,000 & 5 & 6 & \textbf{7} & 8 & 67,7\,\% \\
10\,000 & 10 & 5 & \textbf{5} & 6 & 68,0\,\% \\
\midrule
100\,000 & 3 & 11 & \textbf{14} & 16 & 67,0\,\% \\
100\,000 & 5 & 8 & \textbf{9} & 10 & 67,7\,\% \\
100\,000 & 10 & 6 & \textbf{6} & 7 & 67,9\,\% \\
\bottomrule
\end{tabular}
\end{table}

\paragraph{Tỉ lệ lấp đầy khoảng 68\,\%.} Tỉ lệ lấp đầy $= n / (\text{số nút}\times(m-1))$ ổn định quanh 67--68\,\%
với mọi $m$ và $n$, gần với giá trị lý thuyết xấp xỉ $\ln 2 \approx 69\,\%$ đã biết cho B-Tree chèn khóa theo thứ
tự ngẫu nhiên~\cite{yao}. Nhờ vậy số nút tỉ lệ nghịch với $m$: 74\,673 nút ($m=3$) so với 16\,358 nút ($m=10$).

\paragraph{Thời gian giảm ít hơn số nút.} Khi $m$ tăng từ 3 lên 10, thời gian tìm $10^5$ khóa giảm từ khoảng
108\,ms xuống 62\,ms (khoảng 1,7 lần), ít hơn mức giảm 2,3 lần của số nút (Hình~\ref{fig:thoigian}). Nguyên nhân:
mỗi nút lớn hơn cần nhiều phép so sánh hơn khi tìm nhị phân trong nút, và tổng số phép so sánh chuỗi của B-Tree
vẫn là $O(\log_2 n)$ bất kể $m$. Khi dữ liệu nằm hoàn toàn trong bộ nhớ, lợi ích về thời gian đến chủ yếu từ việc
ít đổi nút hơn (ít truy cập bộ nhớ rải rác). Nếu mỗi nút là một khối trên đĩa, chi phí đọc khối lớn hơn chi phí so
sánh hàng nghìn lần, và số nút truy xuất mới là thước đo quyết định -- đó là bối cảnh B-Tree được thiết
kế~\cite{bayer,comer}.

\paragraph{Thời gian tăng gần tuyến tính theo $n$.} Trên đồ thị log--log, đường thời gian chèn và tìm có độ dốc
xấp xỉ 1 (nhỉnh hơn 1 ở $n$ lớn): tổng chi phí $n$ thao tác là $O(n\log n)$; phần tăng thêm ở $n$ lớn có thể
còn do dữ liệu vượt kích thước bộ nhớ đệm của CPU. Với $n$ nhỏ (dưới 0,5\,ms) thời gian đo bị ảnh hưởng bởi nhiễu của hệ thống
nên các điểm $n = 10^3$ chỉ mang tính tham khảo.

\section{Thời gian các chức năng khác}

Ngoài B-Tree, nhóm đo thêm thời gian của các chức năng trên dữ liệu lớn (số liệu lấy từ các bộ kiểm thử,
biên dịch \code{-O2}):

\begin{table}[htbp]
\centering\small
\caption{Thời gian các chức năng trên dữ liệu lớn}
\label{tab:thoigiankhac}
\begin{tabularx}{\linewidth}{@{}Y p{5.2cm} r@{}}
\toprule
\textbf{Thao tác} & \textbf{Dữ liệu} & \textbf{Thời gian} \\
\midrule
Đọc ba tệp (kiểm tra từng dòng, dựng B-Tree, hàng đợi lô) & $10^4$ mặt hàng, $1{,}9\cdot10^4$ lô, $2{,}8\cdot10^4$ phiếu & 56\,ms \\
Ghi ba tệp & như trên & 38\,ms \\
Tìm theo tên + cảnh báo tồn + cảnh báo HSD + giá trị nhóm + top 10 & $10^5$ mặt hàng, $10^5$ lô & 131\,ms \\
Toàn bộ \code{test\_kho} (gồm 20\,000 thao tác ngẫu nhiên, $m = 3$) & 2\,000 mã vạch & 0,03\,s \\
\bottomrule
\end{tabularx}
\end{table}

Với kho $10^5$ mặt hàng, toàn bộ báo cáo chạy dưới 0,2 giây, đủ nhanh cho thao tác tương tác trên menu.
